% !TeX root = ../../thesis.tex
\markboth{\uppercase{Abstract}}{\uppercase{Abstract}}
\chapter{Abstract}                                 \label{ch:abstract}
\markright{\uppercase{Abstract}}
Translational regulation through synonymous codon usage has recently been shown to play an important role in health and disease. Modified tRNAs are important actors involved in regulating protein expression levels by optimizing the decoding of differentially used codons. Nevertheless their contributions in protein synthesis dynamics remains unclear. Agent-Based Models (ABMs), particularly the Totally Asymmetric Simple Exclusion Process (TASEP), have been employed to quantify ribosomal translation rates. Building on this foundation, our work extends ABM and TASEP frameworks to incorporate the effects of tRNA modifications and additional factors impacting translation. We developed a computational stochastic model to quantify protein synthesis rates, utilizing ribosome residence times at each codon, based on transcript codon sequences and relative transcript abundances. The model integrates codon-specific rates of tRNA accommodation (both modified and unmodified), peptide bond formation, and translocation.

Key features include modulation of elongation rates by charged amino acids within the ribosomal exit tunnel, proline transpeptidation at the peptidyl transferase center, and translational slowing due to mRNA secondary structures. Our model enables the comparison of relative protein expression levels, polysome profiling, and Ribo-Seq data across various scenarios, accounting for ribosome pool availability and initiation rates.

Ultimately, this model aims to elucidate how codon usage and tRNA modifications interact dynamically to influence protein synthesis. By systematically disentangling the contributing factors, we seek to evaluate their sensitivity not only on protein output but also on ribosome density and polysome profiles. %The ultimate purpose of the work is to understand how the translational control through codon usage and tRNA tuning influence the proteome, local protein structure possibly linked to changes in functional specific activity, aggregation, ribophagy, autophagy or proteasomal degradation. The results of this work will possibly help develop targeted therapies against cancers (melanoma and lung cancers).


\instructionsabstract


%%%%%%%%%%%%%%%%%%%%%%%%%%%%%%%%%%%%%%%%%%%%%%%%%%
% Keep the following \cleardoublepage at the end of this file,
% otherwise \includeonly includes empty pages.
\cleardoublepage

% vim: tw=70 nocindent expandtab foldmethod=marker foldmarker={{{}{,}{}}}
